% ------------------------------------------------------------------
%  CHAPTER 2 – LITERATURE SURVEY
% ------------------------------------------------------------------
\chapter{Literature Survey}
 
This chapter provides a comprehensive review of recent advancements in
Retrieval-Augmented Generation (RAG) and intelligent agent frameworks.
The field has progressed considerably from static knowledge retrieval toward dynamic,
corrective, and agentic systems that aim to minimise hallucinations and improve reasoning
depth across a variety of knowledge-intensive tasks.
The survey is structured to trace this progression from foundational contributions
through to the most recent advances in agentic orchestration, memory management,
hallucination mitigation, knowledge-graph integration, and hierarchical retrieval,
culminating in an identification of the research gaps that EvoRAG is specifically
designed to address.

The foundational concepts behind EvoRAG are grounded in the original RAG
framework~\cite{lewis2020rag} and the retrieval-aware reasoning advances exemplified
by ReAct~\cite{yao2022react}.
Work on tool-oriented model orchestration in Toolformer~\cite{schick2023toolformer}
and the self-reflective generation architecture of Self-RAG~\cite{asai2023selfrag}
also directly influence the way EvoRAG integrates external knowledge and internal
consensus validation.
The broader context for these contributions is provided by the foundational study of
large-scale pre-trained language models as general-purpose systems~\cite{bommasani2021foundation},
which established the conceptual framework within which all subsequent retrieval
augmentation research is situated.
 
% ------------------------------------------------------------------
\section{Agentic RAG Frameworks}
% ------------------------------------------------------------------
 
The shift from traditional RAG pipelines to agentic architectures represents one of
the most significant developments in the field over the past two years.
Singh et al.~\cite{singh2025agenticrag} provide a comprehensive survey of Agentic RAG
systems, introducing a principled taxonomy of architectures classified by agent
cardinality, control structure, autonomy, and knowledge representation.
Their analysis demonstrates how embedding autonomous agents into the RAG pipeline
enables dynamic retrieval strategy management, iterative context refinement, and
multi-step reasoning that static, single-pass workflows are fundamentally unable to
achieve.
A key insight from this survey is that no single retrieval granularity is universally
optimal; effective agentic systems must be capable of dynamically selecting between
keyword, semantic, and document-level retrieval depending on the information density
of the query.
 
Complementing this taxonomic perspective, Du et al.~\cite{du2026arag} introduce
A-RAG, a hierarchical retrieval interface framework that exposes keyword search,
semantic search, and chunk-level read tools directly to the language model.
By granting the model the capacity to adaptively navigate across multiple retrieval
granularities, A-RAG consistently outperforms prior approaches while maintaining
comparable or lower retrieved token counts — demonstrating that architectural
flexibility can compensate for brute-force context scaling.
This finding is particularly relevant to EvoRAG's retrieval module, which similarly
distinguishes between coarse topic retrieval and fine-grained passage extraction
depending on the complexity of the query issued by each active persona.
 
More recently, the growing intersection of agentic retrieval and deep reasoning has
been surveyed by Li et al.~\cite{li2025ragreas}, who introduce a unified
Reasoning-Retrieval perspective that categorises methods along two complementary
axes: \emph{Reasoning-Enhanced RAG}, in which advanced reasoning optimises
individual stages of the retrieval pipeline; and \emph{RAG-Enhanced Reasoning}, in
which retrieved knowledge provides the missing factual premises required for complex
multi-step inference.
Their survey further identifies emerging \emph{Synergised RAG-Reasoning} frameworks,
where agentic LLMs iteratively interleave search and reflection to achieve
state-of-the-art performance on knowledge-intensive benchmarks — a paradigm that
directly informs the EvoRAG design.
The survey's emphasis on the mutual dependence of retrieval quality and reasoning
quality underlines why EvoRAG's personas are each assigned a distinct reasoning
style rather than simply sharing a common prompt template.

% ------------------------------------------------------------------
\section{Hierarchical and Graph-Based Retrieval Architectures}
% ------------------------------------------------------------------

A significant limitation of flat chunk-based retrieval is its inability to capture
the hierarchical and relational structure present in real-world document corpora.
Two complementary research directions have emerged to address this: recursive
abstractive retrieval and graph-based retrieval.

Sarthi et al.~\cite{sarthi2024raptor} introduce RAPTOR (Recursive Abstractive
Processing for Tree-Organized Retrieval), a framework that recursively clusters
and summarises retrieved text chunks into progressively higher-level abstractions,
constructing a tree structure over the document corpus.
At query time, retrieval can proceed at any level of the tree, enabling the system
to match coarse thematic queries against high-level summaries while still supporting
fine-grained factual lookups at the leaf level.
RAPTOR demonstrates substantial gains over flat retrieval baselines on multi-step
reasoning tasks, confirming that structural organisation of the retrieval index
yields qualitative improvements in answer coherence.
This hierarchical indexing principle informs EvoRAG's approach to knowledge base
construction, wherein retrieved passages are organised by topic cluster before
being distributed to individual personas.

Edge et al.~\cite{edge2024graphrag} extend the retrieval paradigm further by
proposing GraphRAG, a system that pre-constructs a community-structured knowledge
graph from an input corpus and uses this graph to answer global, query-focused
summarisation tasks.
By traversing the graph at the community level rather than retrieving isolated
passages, GraphRAG captures cross-document relationships and thematic patterns
that are entirely invisible to conventional vector-search retrieval.
The practical implication is that for analytical queries requiring synthesis across
many documents — precisely the class of queries for which EvoRAG is designed —
graph-structured retrieval provides qualitatively richer context than passage-level
retrieval alone.

The challenge of integrating structured knowledge from knowledge graphs with the
generative capabilities of LLMs is surveyed comprehensively by
Pan et al.~\cite{pan2024unifying}, who identify three principal integration
strategies: using LLMs to enhance knowledge graph construction, using knowledge
graphs to constrain LLM generation, and jointly training models over both text and
graph-structured data.
Their analysis highlights that hybrid retrieval — grounding generation in both
unstructured text and structured relational knowledge — consistently outperforms
either modality in isolation, particularly for queries involving entity-level
factual precision.
EvoRAG's architecture acknowledges this finding by designing the retrieval layer
to be modality-agnostic, capable of ingesting both plain text documents and
structured data formats.

% ------------------------------------------------------------------
\section{Hallucination Detection and Reliability in RAG Systems}
% ------------------------------------------------------------------

The tendency of large language models to produce factually incorrect yet
linguistically fluent outputs — a phenomenon widely referred to as hallucination —
represents the primary reliability challenge in deploying LLM-based systems for
high-stakes analytical tasks.
Manakul et al.~\cite{manakul2023selfcheckgpt} introduce SelfCheckGPT, a zero-resource
black-box hallucination detection method that exploits the statistical consistency
of multiple stochastic samples from the same model.
The core insight is that factual claims in a model's output will be consistently
reproduced across independently sampled responses if and only if those claims are
supported by strong parametric knowledge; hallucinated content, by contrast, tends
to vary inconsistently across samples.
SelfCheckGPT achieves state-of-the-art hallucination detection on the WikiBio
benchmark without requiring any reference documents or model internals, making it
directly applicable to black-box deployment scenarios.
This self-consistency principle is directly operationalised in EvoRAG's peer-review
scoring mechanism, where claims that appear consistently across multiple independently
generated persona responses receive higher consensus confidence scores.

Mallen et al.~\cite{mallen2023trust} investigate the conditions under which language
models should be trusted to answer from parametric memory versus those in which
external retrieval is essential.
Their study demonstrates that model confidence is a poor predictor of factual
accuracy for long-tail or temporally sensitive entities, motivating the need for
always-on retrieval rather than retrieval triggered only on model uncertainty signals.
This finding provides empirical justification for EvoRAG's unconditional retrieval
strategy, in which all queries are routed through the knowledge ingestion pipeline
regardless of perceived difficulty, ensuring that no query relies solely on
potentially stale parametric knowledge.

Shi et al.~\cite{shi2023replug} address the complementary problem of augmenting
black-box language models — those whose parameters cannot be modified — with a
retrieval component, through REPLUG (Retrieve and Plug).
REPLUG treats the language model as a scoring function and retrieves documents
whose prepending to the context maximises the model's likelihood on the query,
without any gradient-based fine-tuning of the underlying model.
This plug-and-play design philosophy is directly relevant to EvoRAG's architecture,
which runs retrieval-augmented generation over a locally deployed model whose
weights remain fixed, with adaptation achieved entirely through the evolutionary
persona configuration layer rather than through model retraining.

% ------------------------------------------------------------------
\section{Multi-Agent Collaboration in RAG}
% ------------------------------------------------------------------
 
Multi-agent collaboration has emerged as a particularly productive direction for
handling complex, ambiguous information-seeking tasks that exceed the capacity of any
single-agent pipeline.
Nguyen et al.~\cite{nguyen2025marag} present MA-RAG, a training-free multi-agent
framework that orchestrates four specialised agents — a Planner, a Step Definer, an
Extractor, and a QA Agent — to decompose queries into subtasks including query
disambiguation, evidence extraction, and answer synthesis.
By enabling agents to communicate via chain-of-thought reasoning, MA-RAG achieves
improved accuracy and interpretability on multi-hop and ambiguous question-answering
benchmarks without any model fine-tuning.
The framework's design makes explicit what has long been implicit in multi-step
retrieval pipelines: that different sub-tasks within a complex query require
qualitatively different reasoning strategies, and that separating these strategies
into distinct agents yields more reliable outcomes than conflating them within a
single prompt.
 
The problem of heterogeneous retrieval across structured and unstructured data sources
is addressed by Xiao et al.~\cite{xiao2026massrag}, who propose MASS-RAG, a
Multi-Agent Synthesis approach that assigns distinct roles to separate evidence
summarisation, evidence extraction, and reasoning agents.
By combining their intermediate outputs through a dedicated synthesis stage,
MASS-RAG exposes multiple complementary evidence views before committing to a final
answer, yielding consistent performance improvements over strong single-agent baselines.
Xu et al.~\cite{xu2026anchorrag} address the related challenge of open-world
knowledge graph traversal through AnchorRAG, a framework in which a supervisor agent
formulates iterative retrieval strategies for subordinate retriever agents, which
conduct parallel graph traversal before synthesising the resulting knowledge paths.
This hierarchical delegation strategy improves retrieval robustness and substantially
mitigates the impact of ambiguous or erroneous anchor entities, establishing new
state-of-the-art results on real-world multi-hop reasoning tasks.

The collective insight from these multi-agent frameworks is that the decomposition
of retrieval, reasoning, and synthesis into separate agents controlled by a
coordinator — whether heuristic or learned — is a broadly effective design pattern.
EvoRAG differs from this pattern not in its use of multiple agents but in the
nature of their interaction: rather than specialising each agent by function, EvoRAG
specialises agents by reasoning \emph{style}, allowing the same retrieval context
to be interpreted through distinct analytical lenses before a consensus is formed.

% ------------------------------------------------------------------
\section{Training and Traceability in Agentic RAG}
% ------------------------------------------------------------------
 
End-to-end training of agentic RAG systems using reinforcement learning has shown
strong promise for improving both retrieval reliability and reasoning traceability.
Zheng et al.~\cite{zheng2025deepdx} introduce Deep-DxSearch, an agentic RAG system
trained via reinforcement learning over a corpus of disease profiles, patient records,
and biomedical documents.
By co-optimising retrieval and reasoning through soft verifiable rewards, the system
learns to formulate queries, evaluate evidence, and iteratively refine searches —
directly addressing the disconnect from evidence-based reasoning that is endemic to
static RAG pipelines.
The medical domain context of this work is instructive: it demonstrates that
end-to-end RL training can instil traceable, evidence-grounded reasoning behaviours
even in the absence of direct supervision on intermediate reasoning steps.

The use of human preference signals to align language model behaviour more broadly
was established by Ouyang et al.~\cite{ouyang2022rlhf}, whose InstructGPT system
demonstrated that fine-tuning a language model via reinforcement learning from human
feedback (RLHF) produces outputs that are substantially more aligned with user intent
than those produced by a purely instruction-tuned model of the same scale.
The principle that empirical feedback from evaluators can be used to guide model
behaviour without retraining its core parameters is directly echoed in EvoRAG's
evolutionary engine, which treats user feedback and peer-review win rates as a
lightweight form of preference signal used to promote or demote individual persona
configurations.

% ------------------------------------------------------------------
\section{Data Synthesis for Robust RAG Agents}
% ------------------------------------------------------------------
 
A critical bottleneck in developing robust agentic RAG systems is the scarcity of
high-quality training data that faithfully reflects the retrieval noise encountered in
real-world deployments.
Addressing this gap, RAGShaper~\cite{anonymous2026ragshaper} introduces an automated
data synthesis framework that constructs dense information trees enriched with
adversarial distractor documents.
A constrained navigation strategy forces a teacher agent to actively confront these
distractors during trajectory generation, eliciting demonstrations that explicitly
model error correction and noise rejection.
Models trained on these synthetic trajectories significantly outperform baselines that
train on clean corpora, indicating that controlled adversarial exposure is a viable and
scalable path to improving agent robustness.

The closely related problem of enabling few-shot generalisation in retrieval-augmented
models is addressed by Izacard et al.~\cite{izacard2022atlas}, who introduce Atlas,
a pre-trained retrieval-augmented language model that achieves strong few-shot
performance on knowledge-intensive tasks by jointly training the retriever and the
reader end-to-end.
Atlas demonstrates that a model with access to a large, up-to-date retrieval index
can match or outperform models with orders-of-magnitude more parameters on factual
tasks, provided that the retriever and reader are trained in coordination.
This result motivates EvoRAG's emphasis on tight coupling between the retrieval
pipeline and the persona reasoning layer, rather than treating retrieval as a
preprocessing step decoupled from generation.
 
% ------------------------------------------------------------------
\section{Memory and Long-Term Knowledge in AI Agents}
% ------------------------------------------------------------------
 
Memory management is increasingly recognised as a foundational capability for
intelligent agents operating in complex, multi-session environments.
Hu et al.~\cite{hu2025memory} provide a comprehensive landscape of agent memory
research, proposing a unified taxonomy structured along three axes:
\emph{forms} (token-level, parametric, and latent memory);
\emph{functions} (factual, experiential, and working memory); and
\emph{dynamics} (formation, evolution, and retrieval).
This taxonomy moves beyond the conventional long-term versus short-term dichotomy and
establishes that agent failures are often rooted in deficiencies in memory dynamics
rather than model capacity alone — a finding with direct implications for the EvoRAG
evolutionary learning engine.
 
Complementing this taxonomic foundation, the challenge of \emph{joint} long-term and
short-term memory management has been formalised by
Chen et al.~\cite{chen2026agemem}, who propose AgeMem, a reinforcement learning-based
framework in which an agent learns autonomous policies for promoting, consolidating,
and retrieving information across memory tiers.
Their experiments demonstrate that an RL-trained agent substantially outperforms
heuristic-based memory management strategies on both embodied reasoning and
knowledge-intensive question-answering tasks, underscoring the practical value of
learning-based memory governance.
 
A further dimension of memory research, focused specifically on personalised
conversational continuity, is explored in the Memory OS of AI Agents
framework~\cite{bai2025memos}, which introduces a segmented paging architecture
for dialogue-chain memory.
The system partitions memory into active, inactive, and archival tiers governed by
heat-based mechanisms that model the temporal importance decay of stored episodes.
This operating-system-inspired abstraction enables contextual coherence and
personalised recall across extended, multi-session interactions — a capability that
EvoRAG's user adaptation module is designed to approximate through per-persona
affinity modelling.
The analogy between operating-system memory management and agent knowledge retention
is particularly apt: just as an OS must balance fast access to recently used pages
with long-term persistence of infrequently accessed data, EvoRAG must balance
responsiveness to recent user preferences with the cumulative evolutionary history
that determines which personas have demonstrated sustained analytical value.
 
% ------------------------------------------------------------------
\section{Comparative Evaluation of RAG Approaches}
% ------------------------------------------------------------------
 
As RAG systems grow in architectural complexity, systematic and reproducible evaluation
becomes increasingly critical for guiding practical deployment decisions.
Besrour et al.~\cite{besrour2026agenticvsenh} conduct an experiment-driven comparison
between Enhanced RAG and Agentic RAG across multiple domains, examining the impact of
LLM quality on system robustness and identifying the practical performance trade-offs
between fixed pipeline and agentic orchestration approaches.
Their findings provide actionable guidance for practitioners choosing between the two
paradigms and confirm that agentic orchestration offers measurable advantages on
multi-hop and reasoning-intensive queries, but that this advantage diminishes when
the underlying LLM quality is insufficient to support reliable tool use.
 
Addressing the absence of principled evaluation benchmarks for the RAG field more
broadly, Gan et al.~\cite{gan2025rageval} present a comprehensive survey of RAG
evaluation methods and frameworks, systematically reviewing traditional and emerging
approaches across four dimensions: system performance, factual accuracy, safety, and
computational efficiency.
Their meta-analysis of evaluation practices in high-impact RAG research reveals
persistent inconsistencies in how retrieval and generation quality are measured,
motivating the need for standardised benchmarking protocols of the kind that EvoRAG's
internal peer-review mechanism approximates at the system level.
The study's finding that most existing evaluation frameworks assess retrieval and
generation in isolation — rather than as an integrated system — provides additional
motivation for EvoRAG's end-to-end consensus scoring approach, which measures the
quality of the complete retrieval-reasoning-synthesis pipeline rather than its
individual components.

% ------------------------------------------------------------------
\section{Established Foundations}
% ------------------------------------------------------------------
 
The literature surveyed in the preceding sections builds upon a set of foundational
contributions that established the conceptual and technical basis of the field.
Lewis et al.~\cite{lewis2020rag} introduced the original RAG framework, demonstrating
for the first time that retrieval-augmented generation could substantially improve
performance on knowledge-intensive NLP tasks by grounding generation in non-parametric
external knowledge.
Guu et al.~\cite{guu2020realm} extended this paradigm to the pre-training phase through
REALM, showing that retrieval can be incorporated into language model pre-training
via differentiable document retrieval.
Izacard and Grave~\cite{izacard2021passage} further consolidated the passage-retrieval
paradigm by demonstrating that fusing multiple retrieved passages through a generative
reader yields strong open-domain question-answering performance.
 
On the reasoning side, Wei et al.~\cite{wei2022cot} established chain-of-thought
prompting as the dominant paradigm for eliciting multi-step reasoning in large language
models, providing the conceptual foundation upon which virtually all agentic systems
surveyed above are built.
Yao et al.~\cite{yao2022react} subsequently operationalised this insight in the context
of retrieval by proposing the ReAct framework, which interleaves reasoning traces and
retrieval actions to enable more informed, iterative information-seeking behaviour.
Finally, Asai et al.~\cite{asai2023selfrag} introduced Self-RAG, a self-reflective
generation paradigm in which the model learns to critique both its retrieved evidence
and its generated outputs through specialised reflection tokens — establishing the
principle of internal quality control that directly motivates EvoRAG's consensus-based
peer-review mechanism.
 
% ---------------------------------------------------------------
\section{Research Gaps}

Despite the rapid progression of RAG and agentic AI systems, several important
limitations remain unaddressed in the current body of research, and it is precisely
these gaps that EvoRAG is designed to close.

\begin{itemize}

    \item \textbf{Absence of Parallel Multi-Perspective Reasoning:}
    Most existing RAG and agentic frameworks generate a single final response, even
    when intermediate reasoning steps are involved. Task decomposition across sequential
    agents is not equivalent to the simultaneous generation of competing analytical
    perspectives that are collectively evaluated before a consensus answer is produced.
    No existing system — including MA-RAG~\cite{nguyen2025marag},
    MASS-RAG~\cite{xiao2026massrag}, or AnchorRAG~\cite{xu2026anchorrag} — implements
    this cognitive-diversity model at the intra-query level.

    \item \textbf{Limited Personalisation:}
    Current systems operate in a largely generic fashion and do not adapt their reasoning
    strategies in response to individual user preferences or accumulated feedback over
    time. Memory-based approaches such as AgeMem~\cite{chen2026agemem} and
    MemOS~\cite{bai2025memos} address information retention but do not actively
    modulate reasoning diversity or decision-making styles to match evolving user
    expectations.

    \item \textbf{Absence of Continuous Self-Improvement:}
    Most RAG pipelines remain architecturally static after deployment. Even in advanced
    agentic systems, there is minimal emphasis on dynamically evolving internal reasoning
    components in response to performance metrics or user interaction history.
    While RLHF~\cite{ouyang2022rlhf} and RL-based training~\cite{zheng2025deepdx}
    demonstrate the feasibility of feedback-driven adaptation, neither approach has been
    integrated into a live multi-persona consensus framework that adapts without
    retraining the base model.

    \item \textbf{Dependence on Cloud Infrastructure:}
    A significant proportion of high-performance RAG systems rely on cloud-based APIs
    and external computation resources, introducing concerns related to data privacy,
    regulatory compliance, latency, and operational cost. REPLUG~\cite{shi2023replug}
    demonstrates that black-box augmentation is feasible, but the referenced system
    still depends on externally hosted model endpoints. Locally deployed alternatives
    with comparable analytical depth remain largely unexplored in existing research.

    \item \textbf{Insufficient Handling of Conflicting Information:}
    Current retrieval systems optimise for relevance but frequently fail to explicitly
    resolve contradictory or conflicting information retrieved from multiple sources.
    Even hallucination detection methods such as SelfCheckGPT~\cite{manakul2023selfcheckgpt}
    operate post-hoc on a single generation rather than proactively resolving conflicts
    across multiple retrieved evidence threads. Structured comparison and
    consensus-building mechanisms within RAG pipelines remain underexplored.

    \item \textbf{Static Knowledge Interaction:}
    Even in adaptive retrieval systems such as RAPTOR~\cite{sarthi2024raptor} and
    GraphRAG~\cite{edge2024graphrag}, interaction with retrieved knowledge remains
    largely linear and does not incorporate iterative debate, critique, or mutual
    validation among multiple reasoning agents. The retrieved content informs a single
    generation trajectory rather than seeding a structured deliberation process.

\end{itemize}

These gaps collectively highlight the need for a system that integrates multi-perspective
reasoning, adaptive learning, and privacy-preserving local deployment within a unified
framework. EvoRAG addresses these limitations through a multi-persona reasoning engine,
a consensus-based peer-review mechanism, and an evolutionary learning layer that
continuously refines system performance based on empirical outcome data.